\documentclass[12pt]{article}
\usepackage[utf8]{vietnam}
\usepackage[vietnamese=nohyphenation]{hyphsubst}
\usepackage[vietnamese]{babel}
\usepackage[table,dvipsnames]{xcolor} % Required for row coloring
\usepackage{booktabs} % Required for nicer horizontal lines
\usepackage{ragged2e}
\usepackage{xcolor}
\usepackage{caption}
\usepackage{subcaption}
\usepackage{float}
\usepackage{amsmath,amssymb,amsfonts}
\usepackage{graphicx}
\usepackage{adjustbox}
\usepackage{multicol}
\usepackage[inline,shortlabels]{enumitem}
\usepackage{forest}
\usepackage{setspace}
\usepackage{longtable}
\usepackage{bm} 
\usepackage[inkscapelatex=false]{svg}
\usepackage{titling}
\usepackage{tabularx}
\usepackage{makecell} 
\usepackage[T1]{fontenc}
\usepackage{lmodern,mathrsfs}
% \usepackage[draft]{hyperref}
\usepackage{hyperref}
\usepackage{xparse}
\usepackage[most]{tcolorbox}

%  biblatex (thay cho IEEEtran.bst) ------------------------------------
\usepackage[
    backend=bibtex,
    style=ieee,
    maxnames=1,     % chỉ hiển thị 1 tên tác giả, phần còn lại thành "et al."
    minnames=1,
    uniquelist=false,
    hyperref=false
]{biblatex}

\addbibresource{references.bib}
\DefineBibliographyStrings{english}{
  andothers = {et\addabbrvspace al\adddot}
}
\AtBeginBibliography{\selectlanguage{english}}

% -----------------------------------------------------------------
\AtBeginDocument{\renewcommand\refname{}}

\usepackage[toc,page]{appendix}
\renewcommand{\appendixtocname}{Phụ lục}      
\renewcommand{\appendixpagename}{Phụ lục}

\usepackage{xurl}
\setlength{\headheight}{14.5pt}
\setlength{\topmargin}{-.5in} \setlength{\textheight}{9.25in}
\setlength{\oddsidemargin}{-0.1in} \setlength{\textwidth}{6.8in}

\usepackage{tocloft}
\renewcommand{\cfttoctitlefont}{\Huge\bfseries}
\setlength{\cftaftertoctitleskip}{2em}
\setlength{\cftbeforesubsecskip}{7pt} % giãn cách subitem mục lục
\renewcommand{\cftsecleader}{\cftdotfill{\cftdotsep}}
\setlength{\cftsecnumwidth}{3em}
\renewcommand{\cftsecaftersnum}{\quad}
\renewcommand{\cftsubsecleader}{\cftdotfill{\cftdotsep}}
\setlength{\cftsubsecnumwidth}{4em}
\renewcommand{\cftsubsecaftersnum}{\quad}
\renewcommand{\cftsubsubsecleader}{\cftdotfill{\cftdotsep}}
\setlength   {\cftsubsubsecnumwidth}{5em}  
\renewcommand{\cftsubsubsecaftersnum}{\quad}

\PassOptionsToPackage{vietnamese.licr}{babel}
\usepackage{fontawesome5}
\usepackage{vipythonhighlight} 
\newcounter{mycounter}
\newcommand\showmycounter{\stepcounter{mycounter}\themycounter}
\usepackage{lipsum}
\newcommand\showlips{\stepcounter{mycounter}\lipsum[\value{mycounter}]}

%--------------------------------------------------------------------
\usepackage{framed}
\usepackage{fancyhdr}

\usepackage{listings}	 % Sử dụng gói lệnh listings đề chèn code
\usepackage{tvietlistings}		    % Sử dụng Tiếng Việt trong gói listings

% Color --------------------------------------------------------------------
\usepackage{epigraph}
\usepackage{titlesec}

\definecolor{codebackground}{RGB}{250, 250, 250}
\definecolor{codecomment}{RGB}{34, 139, 34}
\definecolor{codekeyword}{RGB}{0, 0, 255}
\definecolor{codestring}{RGB}{255, 0, 0}
\definecolor{codeidentifier}{RGB}{0, 0, 0}
\definecolor{codelinenumbers}{RGB}{90, 90, 90}

\definecolor{codegreen}{rgb}{0,0.6,0}
\definecolor{codegray}{rgb}{0.5,0.5,0.5}
\definecolor{codepurple}{rgb}{0.58,0,0.82}
\definecolor{backcolour}{rgb}{0.95,0.95,0.92}

\definecolor{aivnGreen}{RGB}{52,79,30}  
\definecolor{captiongray}{gray}{0.15} 
\definecolor{captionbg}{RGB}{250,250,250}

% Table
\newtcolorbox{captionbelowboxaivn}{
    colback=aivnGreen,
    colframe=aivnGreen,
    boxrule=0pt,
    arc=4pt,
    left=4pt,
    right=4pt,
    top=2pt,
    bottom=2pt,
    width=\linewidth,
    boxsep=4pt,
    halign=center,
    % fontupper=\bfseries\color{white}
    fontupper=\color{white}
}

% Figure
\newtcolorbox{figurecardcaptionaivn}{
  colback=aivnGreen,        % Màu nền (background color)
  colframe=aivnGreen,       % Màu viền box (border color)
  boxrule=0pt,              % Độ dày viền box (0pt = không viền)
  arc=2pt,                  % Độ bo góc (2pt là bo nhẹ, tăng để bo tròn hơn)
  top=3pt,                  % Padding phía trên nội dung
  bottom=3pt,               % Padding phía dưới nội dung
  left=10pt,                % Padding bên trái
  right=10pt,               % Padding bên phải
  width=\linewidth,         % Chiều rộng bằng toàn bộ dòng
  fontupper=\small\fontsize{11pt}{13pt}\selectfont\color{white} 
}

% Tạo counter riêng
\newcounter{hinh}
% Định nghĩa macro caption Hình
\newcommand{\capfig}[1]{%
  \refstepcounter{hinh}%
  % \textbf{Hình~\thehinh:}~#1
    Hình~\thehinh:~#1
}

\newtcolorbox{figurecaptionmodern}{
  colback=captionbg,
  colframe=aivnGreen,
  boxrule=1.5pt,
  bottomrule=2pt,  
  toptitle=0pt,
  bottomtitle=0pt,
  left=6pt,
  right=6pt,
  top=4pt,
  bottom=2pt,
  width=\linewidth,
  fontupper=\small\fontsize{11pt}{13pt}\selectfont\color{black},
  halign=center
}

% Code Block -----------------------------------------------------------------
\newtcolorbox{aivncodebox}[1][]{%
  enhanced,
  colback=gray!10,
  colframe=aivnGreen,
  coltitle=white,
  colbacktitle=aivnGreen,
  title=#1,
  fonttitle=\bfseries\sffamily,
  listing only,
  breakable,
  listing engine=listings,
  listing options={
    language=Python,
    basicstyle=\ttfamily\footnotesize,
    commentstyle=\color{ForestGreen},
    keywordstyle=\color{blue},
    stringstyle=\color{red!70!black},
    numbers=left,
    numberstyle=\tiny\color{gray},
    numbersep=5pt,
    breaklines=true,
    showstringspaces=false,
    xleftmargin=1.5em,
    framexleftmargin=1.5em
  }
}

% --- Concept explanation box ---
\newtcolorbox{conceptbox}[2][]{
    enhanced, colback=green!6!white, colframe=aivnGreen, boxrule=1pt, arc=4pt,
    fonttitle=\bfseries, attach boxed title to top left={yshift=-2mm, xshift=3mm},
    colbacktitle=aivnGreen, coltitle=white, title=#2, #1
}

% --- Guideline/Tip box ---
\newtcolorbox{guidelinebox}[2][]{
    enhanced, colback=yellow!10!white, colframe=orange!90!black, boxrule=1pt, arc=4pt,
    fonttitle=\bfseries, attach boxed title to top left={yshift=-2mm, xshift=3mm},
    colbacktitle=orange!90!black, coltitle=white,
    title=\faLightbulb\hspace{1mm} #2, #1
}

% --- Note/Info box ---
\newtcolorbox{notebox}[1][]{
  enhanced,  colback=yellow!10, colframe=orange!80, boxrule=1pt, 
  arc=4pt, fontupper=\small, attach boxed title to top left={yshift=-2mm, xshift=3mm},
  colbacktitle=orange!80,  coltitle=white, fonttitle=\bfseries,  
  title=\faInfoCircle~\textbf{Lưu ý},  #1
}

% Last argument allows more tcolorbox options to be added

%--------------------------------------------------------------------
\lstdefinestyle{mystyle}{
    backgroundcolor=\color{codebackground},
    commentstyle=\color{codecomment},
    keywordstyle=\color{codekeyword},
    stringstyle=\color{codestring},
    identifierstyle=\color{codeidentifier},
    basicstyle=\ttfamily\footnotesize,
    breakatwhitespace=false,
    breaklines=true,
    captionpos=b,
    keepspaces=true,
    numbers=left,
    numberstyle=\tiny\color{codelinenumbers}, 
    numbersep=5pt,
    showspaces=false,
    showstringspaces=false,
    showtabs=false,
    tabsize=2
}

\lstset{style=mystyle}

\hypersetup{
    colorlinks=true,
    linkcolor=black,
    filecolor=magenta,    
    citecolor=black,
    urlcolor=black,
    pdftitle={Overleaf Example},
    pdfpagemode=FullScreen,
}

%--------------------------------------------------------------------
\newcounter{choice}
\renewcommand\thechoice{($\alph{choice}$)}
\newcommand\choicelabel{\thechoice}

%--------------------------------------------------------------------
\newenvironment{choices}%
  {\list{\choicelabel}%
     {\usecounter{choice}\def\makelabel##1{\hss\llap{##1}}%
       \settowidth{\leftmargin}{W.\hskip\labelsep\hskip 0 em}%
       \def\choice{%
         \item
       } % choice
       \labelwidth\leftmargin\advance\labelwidth-\labelsep
       \topsep=0pt
       \partopsep=0pt
     }%
  }%
  {\endlist}
  
%--------------------------------------------------------------------
\def\CChoice{%
      \choice
      \addanswer{\theenumi}{\thechoice}%
    }
    \let\CChoice\CChoice
%    \par % Uncomment this to have choices always start a new line
   % \let\par\@empty
    % If we're continuing the paragraph containing the question,
    % then leave a bit of space before the first choice:
    \ifvmode\else\enskip\fi
    \ignorespaces
  %
  {}
\makeatother
\newbox\allanswers
\setbox\allanswers=\hbox{}
\newcommand{\addanswer}[2]{%
  \global\setbox\allanswers=\hbox{\unhbox\allanswers \quad #1.~#2}%
}
\newcommand{\showanswers}{%
  \vfill
  \begin{center}
    Answers
  \end{center}
  \unhbox\allanswers
}

\makeatletter
\newenvironment{multichoices}[1][1]{%
\begin{multicols}{#1}}{%
\end{multicols}}
\makeatother

\lstset{style=mystyle}
\pagestyle{fancy}

\DeclareMathOperator*{\argmax}{arg\,max}
\DeclareMathOperator*{\argmin}{arg\,min}
\DeclareMathOperator{\E}{\mathbb{E}}

% Outline --------------------------------------------------------------------
\renewcommand{\thesection}{\Roman{section}.} 
\renewcommand{\thesubsection}{\thesection\arabic{subsection}.}
\renewcommand{\thesubsubsection}{\thesubsection\arabic{subsubsection}.} 

\titleformat{\section}{\normalfont\Huge\bfseries}{\thesection}{1em}{}
\titleformat{\subsection}{\normalfont\Large\bfseries}{\thesubsection}{1em}{}
\titleformat{\subsubsection}{\normalfont\normalsize\bfseries}{\thesubsubsection}{1em}{}


% Logo --------------------------------------------------------------------
\setlength{\droptitle}{-8em}
\pretitle{%
    \begin{flushleft}
        \includegraphics[height=2.0cm]{Figures/logo.pdf}\par
    \end{flushleft}
}
\posttitle{}

% Title --------------------------------------------------------------------
\title{%
  \begin{center}
        \Large {AI VIET NAM -- AI COURSE 2025}\\[.5ex]  
        \Huge \textbf{Template}\\[.5ex]
  \end{center}
}

% Author --------------------------------------------------------------------
\posttitle{
    \author{
        \begin{minipage}{\textwidth}
        \centering
        {
          Yen-Linh~Vu,\quad
          Anh-Khoi~Nguyen,\quad
          Truong-Binh~Duong,\quad
          Quang-Vinh~Dinh
        }\\
      \end{minipage}
    }
}
% \setlength{\drop}{-8em}
\date{}

% Header --------------------------------------------------------------------
\pagestyle{fancy}
\fancyhf{}
\lhead{\href{https://www.facebook.com/aivietnam.edu.vn}{\textcolor{aivnGreen}{\bfseries AI VIET NAM (AIO2025)}}}
\rhead{\href{https://aivietnam.edu.vn/}{\textcolor{aivnGreen}{\bfseries aivietnam.edu.vn}}}
\fancyfoot[C]{\thepage}
\renewcommand{\headrulewidth}{1.0pt}
\renewcommand{\headrule}{{\color{aivnGreen}\hrule width\headwidth height\headrulewidth \vskip-\headrulewidth}}

% Document -------------------------------------------------------------------
\begin{document}
\maketitle
\vspace{-3.5em}
\section{Lý thuyết}
\label{section:theory}

Đây là nơi trình bày các khái niệm lý thuyết, dẫn nhập.

\subsection{Ví dụ về các Hộp thông tin}

\begin{conceptbox}{Đây là Hộp Khái niệm \faRandom}
    \textbf{Ý tưởng:} Sử dụng hộp này để định nghĩa hoặc giải thích một khái niệm quan trọng.
\end{conceptbox}

\begin{guidelinebox}{Đây là Hộp thông tin thêm/ mẹo/ tip \faLightbulb}
    Sử dụng hộp này để đưa ra các mẹo, hướng dẫn hoặc các quy tắc thực hành.
    \begin{itemize}
        \item Cho bài toán \textbf{phân loại}: $m \approx \sqrt{n}$.
        \item Cho bài toán \textbf{hồi quy}: $m \approx n/3$.
    \end{itemize}
\end{guidelinebox}

\begin{notebox}
    Đây là hộp lưu ý, dùng để nhấn mạnh một thông tin quan trọng hoặc một cảnh báo cho người đọc.
\end{notebox}


\subsection{Ví dụ về Bảng Ký hiệu Toán học}

\begin{tcolorbox}[
    colback=white, colframe=aivnGreen, boxrule=1pt,
    title=\textbf{Bảng Ký hiệu Toán học}, fonttitle=\bfseries\color{white},
    halign title=center, colbacktitle=aivnGreen, arc=5pt
]
\renewcommand{\arraystretch}{1.5}
\rowcolors{2}{white}{green!12!white}
\begin{tabular}{>{\centering\arraybackslash}p{2.5cm} p{12cm}}
    \textbf{Ký hiệu} & \textbf{Ý nghĩa} \\
    \midrule
    $B$ & Siêu tham số, là tổng số cây quyết định trong rừng. \\
    $N$ & Tổng số điểm dữ liệu trong tập huấn luyện gốc. \\
    $\mathcal{D}$ & Tập dữ liệu huấn luyện gốc, $\mathcal{D} = \{(\mathbf{p}_1, y_1), \dots, (\mathbf{p}_N, y_N)\}$. \\
\end{tabular}
\end{tcolorbox}


\subsection{Ví dụ về Layout 2 cột}

\begin{figure}[H]
\begin{minipage}[c]{0.45\textwidth}
    \centering
    \begin{forest}
        for tree={
            grow=south,
            parent anchor=south,
            child anchor=north,
            l sep=8mm,
            s sep=2mm,
            rounded corners,
            draw,
            align=center,
            font=\small
        }
        [``Thời tiết?'', fill=blue!20
            [``Nắng''
                [``Chơi'', fill=green!40]
            ]
            [``U ám''
                [``Chơi'', fill=green!40]
            ]
            [``Mưa''
                [``Nghỉ'', fill=red!30]
            ]
        ]
    \end{forest}
    \captionof{figure}{Cây quyết định minh họa.}
    \label{fig:dt_illustration}
\end{minipage}%
\begin{minipage}[c]{0.55\textwidth}
    \centering
    \begin{tcolorbox}[
        colback=white, colframe=blue!50, boxrule=1pt, arc=2pt,
        title=\textbf{\faInfoCircle\ Chú thích}, fonttitle=\small\bfseries,
        halign title=center, top=2pt, bottom=2pt
    ]
    \small
    \begin{itemize}[leftmargin=*, topsep=3pt, itemsep=1pt, parsep=0pt]
        \item \textbf{Nút gốc (Root Node):} Nút bắt đầu của cây (``Thời tiết?'').
        \item \textbf{Nhánh (Branch):} Kết quả của một phép thử, nối các nút lại với nhau (``Nắng'', ``U ám'', ``Mưa'').
        \item \textbf{Nút lá (Leaf Node):} Nút cuối cùng, đưa ra dự đoán (``Chơi'', ``Nghỉ'').
    \end{itemize}
    \end{tcolorbox}
\end{minipage}
\end{figure}
\hline
\begin{figure}[H]
\centering
\begin{minipage}{0.48\textwidth}
\centering
\renewcommand{\arraystretch}{1.3}
\setlength{\tabcolsep}{12pt}
\begin{tabular}{|c|c|}
\hline
\textbf{Điểm thi thử} & \textbf{Điểm tốt nghiệp} \\
\hline
14.0   & 12.0   \\ \hline
15.0   & 13.0   \\ \hline
16.5   & 14.0   \\ \hline
17.25  & 14.5   \\ \hline
22.0   & 20.0   \\ \hline
23.5   & 20.5   \\ \hline
26.0   & 27.0   \\ \hline
27.5   & 27.5   \\ \hline
29.0   & ?      \\ \hline
\end{tabular}
\captionof{table}{Dữ liệu điểm thi thử và điểm tốt nghiệp.}
\label{tab:exam-graduation}
\end{minipage}
\hfill
\begin{minipage}{0.48\textwidth}
\RaggedRight
% Box 1: Phân tích
\begin{tcolorbox}[
  colback=gray!5!white,
  colframe=aivnGreen,
  fonttitle=\bfseries,
  title=\faInfoCircle\;Phân tích dữ liệu ban đầu,
  boxrule=0.8pt,
  arc=3pt
]
Dữ liệu này chỉ gồm một biến đầu vào duy nhất, do đó đây là bài toán hồi quy đơn biến.
Quan sát bảng trên, ta thấy mối quan hệ giữa điểm thi thử và điểm tốt nghiệp gần như tuyến tính: điểm thi thử tăng thì điểm tốt nghiệp cũng tăng theo.
\end{tcolorbox}

\vspace{0.8em}

% Box 2: Yêu cầu
\begin{tcolorbox}[
  colback=orange!7!white,
  colframe=orange!80!black,
  fonttitle=\bfseries,
  title=\faTasks\;Yêu cầu,
  boxrule=0.8pt,
  arc=3pt
]
Xây dựng \textbf{Regression Tree} với độ sâu bằng 2 (\texttt{depth = 2}) từ dữ liệu trên, tìm ra các ngưỡng phân tách thích hợp trên trục \emph{điểm thi thử}, và sử dụng mô hình để dự đoán điểm tốt nghiệp cho học sinh có điểm thi thử \textbf{29}.
\end{tcolorbox}

\end{minipage}
\end{figure}
% SECTION II: EXERCISES / PRACTICE ===============================================
\section{Bài tập}
\label{section:practice}

Phần này dùng để mô tả bài tập thực hành.

\subsection{Mô tả}
Mô tả bài toán, vấn đề.

\subsection{Ví dụ về Code Block}

\begin{aivncodebox}[student\_management.py]
\begin{python}
# Step 1: Declare variables
student_name = "An"
math_score = 8.5
science_score = 7.0

# Step 2: Print initial information
# Method 1: Using + to concatenate strings
print("Tên: " + student_name + ", Điểm Toán: " + str(math_score) + ", Điểm Khoa học: " + str(science_score))

# Method 2: Using f-string
print(f"Tên: {student_name} - Toán: {math_score} - Khoa học: {science_score}")
\end{python}
\end{aivncodebox}



% SECTION III: MULTIPLE CHOICE QUESTIONS =========================================
\newpage
\section{Câu hỏi trắc nghiệm}
\label{section:multiplechoices}

\begin{enumerate}[leftmargin=*]
    \item Câu hỏi ví dụ số 1?
    \begin{enumerate}[label=(\alph*)]
        \item Lựa chọn A.
        \item Lựa chọn B.
        \item Lựa chọn C.
        \item Lựa chọn D.
    \end{enumerate}

    \item Câu hỏi ví dụ số 2?
    \begin{enumerate}[label=(\alph*)]
        \item \pyth{agent.tool}
        \item \pyth{@smolagents.tool}
        \item \pyth{@tool}
        \item \pyth{@custom_tool}
    \end{enumerate}
\end{enumerate}

% SECTION IV: REFERENCES =========================================================
\newpage
\section{Tài liệu tham khảo}
\label{section:references}
\nocite{*} % Cite all entries from the .bib file
\printbibliography

% APPENDIX =======================================================================
\newpage
\begin{appendices}

\section{Nội dung Phụ lục}

\begin{enumerate}
    \item \textbf{Datasets:} Các file dataset được đề cập trong bài có thể được tải tại \href{https://aivietnam.edu.vn/}{đây}.
    \item \textbf{Hint:} Các file code gợi ý có thể được tải tại \href{https://aivietnam.edu.vn/}{đây}.
    \item \textbf{Solution:} Lời giải chi tiết có thể được tải tại \href{https://aivietnam.edu.vn/}{đây}.
    \item \textbf{Rubric:}

\begin{longtable}{|c|p{7.5cm}|p{7cm}|}
    \hline
    \textbf{Mục} & \textbf{Kiến Thức} & \textbf{Đánh Giá} \\\hline
    I. &
    - Kiến thức về... \newline
    - Kiến thức về...
    &
    - Nắm được lý thuyết... \newline
    - Có khả năng triển khai...
    \\\hline
\end{longtable}
\end{enumerate}

\end{appendices}

\end{document}
